\section{Autoregressive inference and KV state}
\label{sec:autoregressive-kv}

\subsection{Prefill and decode}

Autoregressive generation selects a new token from the model's
next-token distribution and extends the prefix before producing another:
\begin{equation}
\label{eq:autoregressive-factorization}
 p(x_1,\ldots,x_T)=\prod_{t=1}^{T}p(x_t\mid x_{<t}).
\end{equation}
\eqcaption{Autoregressive factorization: each token is conditioned on
the preceding tokens.}

The model's logits can be turned into probabilities and a token can be
selected deterministically or by sampling.  Generation stops at a
termination token or another declared limit.  These choices affect the
output, but the important execution dependency is that a sequence's
next token must be known before its following evaluation can proceed.
Different sequences can nevertheless advance together.

\textbf{Prefill} evaluates the prompt.  Its positions can be processed
in parallel within matrix operations because the token IDs are already
known, while the causal mask prevents attention to future positions.
For ordinary generation, the logits at the last prompt position produce
the first output token.  \textbf{Decode} then evaluates newly generated
tokens, conventionally one new position per active sequence per step.
Evaluating that first output token produces the distribution for the
second output token.  Thus producing \(O\) output tokens normally needs
prefill and \(O-1\) subsequent decode evaluations; the final emitted
token need not itself be evaluated unless further generation is needed.
Chunked prompt processing and other execution strategies can change how
these logical phases are scheduled, without changing this dependency.

\subsection{What the cache retains}

At each layer, the input activations generate keys and values using the
layer's trained projections.  A later position's query compares against
the available keys, and the resulting attention coefficients combine
the associated values.  Since earlier positions cannot attend to later
ones under causal attention, their keys and values need not be recomputed
when the sequence is extended.  The \textbf{KV cache} retains these
per-layer values for reuse~\cite{pope2023efficient}.

Prefill establishes entries for the evaluated prompt positions.  Decode
reads the retained history and adds entries for the position being
processed.  Queries and most intermediate activations are temporary;
keys and values persist for as long as the corresponding history is
needed.  The cache avoids repeating computation of the previous
positions' KV state; it does not eliminate reading that state or
computing attention for a new query.  With full-history attention, longer
histories increase both retained payload and the state a query can access.
The logical attention expression need not be implemented by materializing
its whole score matrix; kernels can operate on tiles.

The number of \textbf{KV heads} controls how many distinct key/value
vectors are retained per position.  In multi-head attention (MHA), each
query head has its own KV head.  Multi-query attention (MQA) shares one
KV head across all query heads.  Grouped-query attention (GQA) shares a
KV head within each group of query heads, giving an intermediate KV
count~\cite{ainslie2023gqa}.  For fixed layer count, head dimension,
precision and retained length, fewer KV heads reduce the KV payload
proportionally.  The query-head count can remain larger, so a small KV
payload does not mean that all attention computation shrinks by the same
factor.  The attention arrangement is a property of the model architecture,
not a storage policy chosen freely by an allocator.

\subsection{Logical history and physical storage}

A sequence has a logical set of valid KV entries.  A runtime may store
these in a contiguous buffer or map them to fixed-size physical blocks.
Block allocation accommodates sequences of different lengths without
requiring one contiguous maximum-length buffer for each
sequence~\cite{kwon2023pagedattention}.  The PagedAttention work develops
this mechanism for GPU serving; its distinction between logical tokens
and physical blocks is also useful when describing CPU runtime storage,
without assuming identical implementation or performance.

Four quantities must remain distinct.  \textbf{Populated entries} hold
valid state for evaluated positions.  \textbf{Occupied blocks} include
storage assigned to these entries and any unused space within their
blocks.  \textbf{Reserved pool capacity} is the storage set aside for
KV management, including blocks not currently assigned to requests.
\textbf{Resident memory} is the portion currently backed by physical
pages in RAM.  A reserved virtual allocation, a committed pool and its
resident pages can have different sizes.  Lazy page population, eager
initialization and operating-system behavior affect their relationship.

For example, with 16 positions per block, a 17-position sequence occupies
two blocks: 32 positions of block capacity but only 17 populated entries.
A pool of 100 such blocks reserves 1{,}600 positions of capacity even if
only those two blocks are occupied.  Completing the request releases its
logical ownership, but a runtime can keep the pool allocated for later
requests.  Process memory need not drop when the request finishes.
\figref{fig:inference-memory-lifetimes} shows these distinct lifetimes.

\begin{figure}[htbp]
\centering
% Storyboard: one request, shared weights, KV slots inside a reusable pool.
% Filled cells represent valid state; clearing ownership need not erase bytes.
\begingroup
\resizebox{\textwidth}{!}{%
\begin{tikzpicture}[
 font=\small\sffamily,text=black!80,
 weight/.style={draw=figcol,fill=figcol!8,line width=0.65pt},
 temp/.style={draw=tabline,fill=neutralcol,line width=0.65pt},
 kv/.style={draw=ddrcol,fill=ddrcol!25,line width=0.65pt},
 free/.style={draw=ddrcol!45,fill=white,line width=0.5pt},
 flow/.style={-{Stealth[length=4pt,width=3pt]},draw=black!65,line width=0.8pt},
 note/.style={font=\footnotesize\sffamily,align=center},
 stage/.style={font=\small\sffamily\bfseries,align=center},
 rowlabel/.style={anchor=east,align=right,text width=2.0cm},
 token/.style={temp,rounded corners=1pt,minimum width=0.42cm,
   minimum height=0.38cm,inner sep=1pt,font=\footnotesize\sffamily},
 pool/.style={draw=ddrcol,dashed,rounded corners=3pt,line width=0.7pt}
]

% The request path: files, prompt, parallel prompt work, serial steps, text.
\node[stage] at (0,1.7) {Load};
\node[stage] at (3,1.7) {Prompt};
\node[stage] at (6,1.7) {Prefill};
\node[stage] at (9,1.7) {Decode};
\node[stage] at (12,1.7) {Complete};

% Document icon used for model loading.
\draw[weight] (-0.48,-0.25)--(0.48,-0.25)--(0.48,0.57)
 --(0.23,0.82)--(-0.48,0.82)--cycle;
\draw[draw=figcol] (0.23,0.82)--(0.23,0.57)--(0.48,0.57);
\foreach \x in {0,1,2} \foreach \y in {0,1}
 \draw[draw=figcol!65,fill=white]
  (-0.31+\x*0.21,0.05+\y*0.21) rectangle ++(0.16,0.16);
\node[note] at (0,-0.65) {Model files};

% Text enters as a prompt; token identifiers are shown below it.
\draw[temp,rounded corners=3pt] (2.12,-0.12) rectangle (3.88,0.85);
\draw[temp] (2.38,-0.12)--(2.38,-0.3)--(2.64,-0.12);
\node[note] at (3,0.4) {Explain\\caches};
\foreach \j/\dx in {1/-0.55,2/0,3/0.55}
 \node[token] at (3+\dx,-0.65) {$p_{\j}$};
\draw[flow] (3,-0.13)--(3,-0.43);
\node[note,text=black!60] at (3,-1.13) {Token IDs};

% Small processor icons: several known prompt positions, then one new position.
\foreach \cx in {6,9} {
 \foreach \dy in {-0.15,0.15,0.45,0.75} {
  \draw[draw=black!45] (\cx-0.6,\dy)--(\cx-0.77,\dy);
  \draw[draw=black!45] (\cx+0.6,\dy)--(\cx+0.77,\dy);
 }
 \draw[temp,rounded corners=2pt] (\cx-0.6,-0.32) rectangle (\cx+0.6,0.93);
}
\foreach \dx in {-0.35,0,0.35}
 \draw[temp,fill=white] (6+\dx-0.13,0.15) rectangle ++(0.26,0.4);
\node[token] at (9,0.35) {$g_t$};
\node[note] at (6,-0.75) {Prompt positions};
\node[note] at (9,-0.75) {One new position};
\draw[flow] (9.6,0.95)--(9.6,1.25)--(8.4,1.25)--(8.4,0.95);

% Emission follows prefill as well as decode; the loop evaluates emitted tokens.
\draw[flow] (0.7,0.35)--(1.93,0.35);
\draw[flow] (4.03,0.35)--(5.2,0.35);
\draw[flow] (6.8,0.35)--(8.2,0.35);
\node[token] at (7.5,0.35) {$g_1$};
\node[note,text=black!60] at (7.5,-0.3) {First token};
\draw[flow] (9.8,0.35)--(10.92,0.35);
\draw[temp,rounded corners=3pt] (11.1,-0.12) rectangle (12.9,0.85);
\draw[temp] (12.45,-0.12)--(12.45,-0.3)--(12.19,-0.12);
\node[note] at (12,0.4) {Caches\\reuse data.};
\foreach \j/\dx in {1/-0.55,2/0,3/0.55}
 \node[token] at (12+\dx,-0.65) {$g_{\j}$};
\draw[flow] (12,-0.43)--(12,-0.13);
\node[note,text=black!60] at (12,-1.13) {Tokens to text};

\draw[draw=tabline,dashed] (-2.55,-1.55)--(13.35,-1.55);

% The same trained matrices remain available for the engine's lifetime.
\node[rowlabel] at (-1.45,-2.55) {Model\\weights};
\draw[draw=figcol!40,line width=0.7pt] (-0.8,-2.4)--(12.8,-2.4);
\foreach \cx in {0,6,9,12} {
 \draw[weight] (\cx-0.43,-2.75) rectangle (\cx+0.43,-2.05);
 \foreach \dx in {-0.14,0.14}
  \draw[draw=figcol!50] (\cx+\dx,-2.75)--(\cx+\dx,-2.05);
 \draw[draw=figcol!50] (\cx-0.43,-2.4)--(\cx+0.43,-2.4);
}
\node[note,fill=white,inner sep=3pt] at (3,-2.4) {Same weights};
\node[note,text=black!60] at (9,-3.2) {Available across requests};

% Three snapshots of one fixed-capacity pool, each with three two-position blocks.
% A slot comprises the K and V cells in one column; widths are illustrative.
\node[rowlabel] at (-1.45,-4.7) {KV cache\\and pool};
\node[note,text width=2.1cm,text=black!60] at (1.45,-4.7)
 {Allocation time\\depends on engine};
\foreach \cx/\used in {6/3,9/5,12/0} {
 \draw[pool] (\cx-1.05,-5.28) rectangle (\cx+1.12,-4.07);
 \node[note,anchor=east] at (\cx-1.2,-4.46) {$K$};
 \node[note,anchor=east] at (\cx-1.2,-4.91) {$V$};
 \foreach \block in {0,1,2} {
  \pgfmathtruncatemacro{\blockstart}{2*\block}
  \ifnum\blockstart<\used
   \draw[draw=ddrcol,line width=0.75pt,rounded corners=1pt]
    (\cx-0.91+\block*0.67,-5.12) rectangle ++(0.61,0.91);
  \else
   \draw[draw=ddrcol!45,dashed,line width=0.6pt,rounded corners=1pt]
    (\cx-0.91+\block*0.67,-5.12) rectangle ++(0.61,0.91);
  \fi
  \foreach \slot in {0,1} {
   \pgfmathtruncatemacro{\entry}{2*\block+\slot}
   \foreach \dy in {-4.6,-5.03} {
    \ifnum\entry<\used
     \draw[kv] (\cx-0.85+\block*0.67+\slot*0.26,\dy)
      rectangle ++(0.23,0.31);
    \else
     \draw[free] (\cx-0.85+\block*0.67+\slot*0.26,\dy)
      rectangle ++(0.23,0.31);
    \fi
   }
  }
 }
}
\node[note] at (6,-5.75) {Prompt KV};
\node[note] at (9,-5.75) {Growing history};
\node[note] at (12,-5.75) {Slots reusable};
\node[note,text=black!60] at (12,-6.18) {Ownership released};

% Stage-local data: compact input buffers, activation arrays and reusable buffers.
\node[rowlabel] at (-1.45,-7.12) {Temporary\\data};
\foreach \cx in {0,3,6,9} {
 \draw[temp,rounded corners=2pt] (\cx-0.65,-7.49) rectangle (\cx+0.65,-6.66);
}
\foreach \cx in {0,6} \foreach \dx in {-0.4,0,0.4} \foreach \dy in {-7.3,-6.98}
 \draw[draw=tabline,fill=white] (\cx+\dx-0.13,\dy) rectangle ++(0.26,0.24);
\foreach \dx in {-0.4,0,0.4} {
 \draw[draw=tabline,fill=white] (3+\dx-0.13,-7.23) rectangle ++(0.26,0.32);
 \draw[draw=tabline,fill=white] (9+\dx-0.13,-7.23) rectangle ++(0.26,0.32);
}
\draw[draw=tabline,dashed,rounded corners=2pt]
 (11.35,-7.49) rectangle (12.65,-6.66);
\draw[flow] (11.62,-7.05) arc[start angle=180,end angle=-120,radius=0.3cm];
\node[note] at (0,-7.96) {Load buffers};
\node[note] at (3,-7.96) {Input IDs};
\node[note] at (6,-7.96) {Prompt activations};
\node[note] at (9,-7.96) {Step activations};
\node[note] at (12,-7.96) {Buffers reusable};

% Keys distinguish populated KV data from capacity without colour-name labels.
\node[kv,minimum width=0.25cm,minimum height=0.22cm,inner sep=0pt]
 at (0,-8.9) {};
\node[note,anchor=west] at (0.27,-8.9) {Valid K/V};
\node[free,minimum width=0.25cm,minimum height=0.22cm,inner sep=0pt]
 at (4.2,-8.9) {};
\node[note,anchor=west] at (4.47,-8.9) {Unused capacity};
\draw[pool] (8.48,-9.05) rectangle (8.78,-8.75);
\node[note,anchor=west] at (9.02,-8.9) {Retained pool};
\end{tikzpicture}%
}
\endgroup

\caption{Request stages and reusable storage in a persistent engine.
Prefill evaluates prompt tokens \(p_i\) and emits \(g_1\); decode
evaluates generated tokens sequentially.  KV state grows from three to
five evaluated positions within a six-position pool.  Paired K/V cells
hold valid entries; solid block outlines mark assigned blocks.
Completing this unshared request releases its slots while the pool can
remain allocated.  Empty cells indicate unused capacity, not erased
bytes.  Weights persist; temporary buffers may be reused.  Token examples
and dimensions are schematic, and pool allocation time and physical
residency depend on the runtime and memory policy.  The final emitted
token need not have a KV entry.}
\label{fig:inference-memory-lifetimes}
\end{figure}
\FloatBarrier

Sequences are not always physically independent.  \textbf{Prefix caching}
can reuse compatible KV blocks for a common prefix; reference counting
allows several requests to share these blocks.  Compatibility depends
on the model and the context that produced the state, not merely on
matching a suffix of text.  Shared blocks must remain valid while in
use, and an engine may keep reusable prefix blocks after individual
requests complete~\cite{vllmprefix2026}.  Summing every sequence's logical
length therefore gives an unshared payload estimate, not necessarily
the number of physical bytes occupied.  Sliding-window or other bounded
attention arrangements can also limit retained history.  Capacity models
must state which attention and sharing assumptions they use.
\FloatBarrier
